% ==============================================================================
% Figura: Diagrama de Bloques Atómico — Servo 0° con Reinicio de Clegg y Freno
% Estados: MOTOR_PID_POS / STOP_HOME / Fase F1 de RECOVER_HOME (1 kHz)
% Paleta, estilo y 3 zonas idénticas a Figs. 6, 7, 8 y 10. Solo nombres cortos
% en bloques, realimentación cerrada con posición theta y velocidad dtheta_f1
% en pistas estrictamente separadas.
% Bibliotecas TikZ requeridas: arrows.meta, positioning, fit, calc,
%                              backgrounds, babel, bending.
% ==============================================================================
\begin{figure*}[!t]
\centering
% --- Paleta semántica armónica (idéntica a Figs. 6, 7, 8 y 10) ---
\definecolor{blkInk}{HTML}{1E293B}
\definecolor{blkPendEdge}{HTML}{1B4965}
\definecolor{blkPendFill}{HTML}{EAF4FB}
\definecolor{blkArmEdge}{HTML}{0F766E}
\definecolor{blkArmFill}{HTML}{E6F7F5}
\definecolor{blkCondEdge}{HTML}{B45309}
\definecolor{blkCondFill}{HTML}{FEF6EB}
\definecolor{blkDistEdge}{HTML}{7E2268}
\definecolor{blkDistFill}{HTML}{FBF0F8}
\definecolor{blkActEdge}{HTML}{A63429}
\definecolor{blkActFill}{HTML}{FDF0EE}
\definecolor{blkPlantEdge}{HTML}{243447}
\definecolor{blkPlantFill}{HTML}{ECEFF4}
\definecolor{blkZoneZ1}{HTML}{33658A}
\definecolor{blkZoneL2}{HTML}{A63429}
\definecolor{blkZoneL3}{HTML}{2A6B43}

% --- Parámetros centralizados de ritmo geométrico (Layout by Construction) ---
\def\scBlkW{2.04cm}
\def\scBlkH{0.84cm}
\def\scColGapX{0.58cm}
\def\scErrTapX{0.82cm}
\def\scErrLeadX{0.58cm}
\def\scSumCGapX{0.92cm}
\def\scCascGapX{0.54cm}
\def\scActGapX{0.54cm}
\def\scSensGapX{0.56cm}
\def\scCleggGapY{0.62cm}
\def\scPropGapY{0.50cm}
\def\scDerivDropY{0.78cm}
\def\scHystDropY{0.66cm}
\def\scHystInOff{0.18cm}
\def\scSelOffX{0.26cm}
\def\scHopRad{0.095cm}
\def\scTrkStepX{0.44cm}

\resizebox{\textwidth}{!}{%
\begin{tikzpicture}[
    >=Stealth,
    font=\sffamily,
    % --- Estilos de bloques atómicos (solo nombre corto, sin parámetros ni fórmulas) ---
    baseblk/.style={
        rectangle,
        rounded corners=3.5pt,
        align=center,
        inner xsep=3.5pt,
        inner ysep=3pt,
        text=blkInk,
        text width=\scBlkW,
        minimum height=\scBlkH,
        font=\small\sffamily\bfseries
    },
    armblk/.style={baseblk, draw=blkArmEdge, line width=1.1pt, fill=blkArmFill},
    intblk/.style={baseblk, draw=blkPendEdge, line width=1.1pt, fill=blkPendFill},
    cleggblk/.style={baseblk, draw=blkDistEdge, line width=1.1pt, fill=blkDistFill},
    condblk/.style={baseblk, draw=blkCondEdge, line width=1.1pt, fill=blkCondFill},
    selblk/.style={baseblk, draw=blkPlantEdge, line width=1.2pt, fill=blkPlantFill},
    actblk/.style={baseblk, draw=blkActEdge, line width=1.1pt, fill=blkActFill, text width=1.94cm},
    plantblk/.style={baseblk, draw=blkPlantEdge, line width=1.35pt, fill=blkPlantFill, text width=1.82cm},
    sensTblk/.style={baseblk, draw=blkArmEdge, line width=1.1pt, fill=blkArmFill, text width=2.14cm},
    sum/.style={
        circle,
        draw=blkInk,
        line width=1.05pt,
        fill=white,
        minimum size=0.58cm,
        inner sep=0pt,
        font=\small\bfseries,
        text=blkInk
    },
    dot/.style={circle, fill=#1, inner sep=1.5pt},
    % --- Flechas canónicas ---
    wire/.style={line width=1.0pt, draw=blkInk},
    sig/.style={-{Stealth[length=4.6pt, width=3.4pt]}, line width=1.0pt, draw=blkInk},
    sigshort/.style={-{Stealth[length=3.9pt, width=2.9pt]}, line width=1.0pt, draw=blkInk},
    fbT/.style={-{Stealth[length=4.4pt, width=3.1pt]}, line width=0.95pt, draw=blkArmEdge},
    lineT/.style={line width=0.95pt, draw=blkArmEdge},
    sigClegg/.style={-{Stealth[length=4.0pt, width=2.9pt]}, line width=0.95pt, densely dashed, draw=blkDistEdge},
    sigBrakeDash/.style={-{Stealth[length=4.0pt, width=2.9pt]}, line width=0.95pt, densely dashed, draw=blkActEdge},
    wireBrakeDash/.style={line width=0.95pt, densely dashed, draw=blkActEdge},
    sigBrakeSolid/.style={-{Stealth[length=4.2pt, width=3.1pt]}, line width=1.05pt, draw=blkActEdge},
    % --- Contenedores y etiquetas (idénticos a Figs. 6, 7, 8 y 10) ---
    zonebox/.style={
        rectangle,
        rounded corners=5pt,
        draw=#1,
        densely dashed,
        line width=0.95pt,
        fill=#1,
        fill opacity=0.035,
        text opacity=1
    },
    zonetitle/.style={
        rectangle,
        rounded corners=3pt,
        draw=#1,
        line width=0.8pt,
        fill=white,
        inner xsep=5.5pt,
        inner ysep=2.2pt,
        font=\small\sffamily\bfseries,
        text=#1
    },
    signlbl/.style={font=\footnotesize\bfseries, text=blkInk, inner sep=1pt},
    wirelbl/.style={font=\footnotesize\sffamily, text=blkInk, inner sep=1.5pt}
]

% ==============================================================================
% ZONA 1: CONTROL SERVO 0° CON REINICIO DE CLEGG Y FRENO — ANCLAJE EN (0,0)
% ==============================================================================
\node[sum] (sig_p) at (0,0) {$\Sigma_p$};

\coordinate[right=\scErrTapX of sig_p.east] (split_e);

\node[intblk, right=\scErrLeadX of split_e] (win_box) {Ventana\\Integral};
\node[intblk, right=\scColGapX of win_box] (int_box) {Integrador};
\node[intblk, right=\scColGapX of int_box] (clamp_box) {Clamp\\Integral};

\node[sum, right=\scSumCGapX of clamp_box] (sig_c) {$\Sigma_c$};

\node[condblk, right=\scCascGapX of sig_c] (sat_box) {Saturaci\'on};
\node[condblk, right=\scCascGapX of sat_box] (slew_box) {Slew-Rate};
\node[selblk,  right=\scCascGapX of slew_box] (sel_box) {Selector\\Freno};

\node[cleggblk, above=\scCleggGapY of int_box] (clegg_box) {Reinicio\\Clegg};
\node[armblk,   above=\scPropGapY  of clegg_box] (prop_box) {Proporcional};

\node[armblk, below=\scDerivDropY of clamp_box] (deriv_box) {Derivativo};
\node[actblk, text width=\scBlkW] (brake_box) at (slew_box |- deriv_box) {Freno\\Activo};

\coordinate[below=\scHystDropY of deriv_box.south] (hyst_top_row);
\node[actblk, text width=\scBlkW, anchor=north] (hyst_box) at (int_box |- hyst_top_row) {Hist\'eresis};

\coordinate (z1_top)   at ([yshift=0.42cm]prop_box.north);
\coordinate (z1_bot)   at ([yshift=-0.32cm]hyst_box.south);
\coordinate (z1_left)  at ([xshift=-0.62cm]sig_p.west);
\coordinate (z1_right) at ([xshift=0.34cm]sel_box.east);

\begin{scope}[on background layer]
    \node[zonebox=blkZoneZ1, inner sep=0pt,
          fit=(z1_top)(z1_bot)(z1_left)(z1_right)] (zone1) {};
\end{scope}
\node[zonetitle=blkZoneZ1, anchor=west] at ([xshift=0.35cm]zone1.north west)
    {Control Servo $0^\circ$ con Reinicio de Clegg (Paralelo + Cascada)};

% ==============================================================================
% ZONA 2: ACTUACIÓN Y MOTOR (CASCADA IDÉNTICA A FIGS. 6, 7, 8 Y 10)
% ==============================================================================
\coordinate (z2_y) at ([yshift=-1.92cm]zone1.south);
\node[actblk, anchor=west] (banda) at ([xshift=3.60cm]zone1.west |- z2_y) {Banda\\Silenciosa};
\node[actblk, right=\scActGapX of banda] (rampa) {Rampa\\L298N};
\node[actblk, right=\scActGapX of rampa] (pwm) {Modulador\\PWM};
\node[actblk, right=\scActGapX of pwm] (puente) {Puente H\\L298N};
\node[plantblk, right=0.72cm of puente] (motor) {Motor\\80:1};

\coordinate (z2_top)   at ([yshift=0.44cm]banda.north);
\coordinate (z2_bot)   at ([yshift=-0.30cm]banda.south);
\coordinate (z2_left)  at ([xshift=-0.30cm]banda.west);
\coordinate (z2_right) at ([xshift=0.32cm]motor.east);

\begin{scope}[on background layer]
    \node[zonebox=blkZoneL2, inner sep=0pt,
          fit=(z2_top)(z2_bot)(z2_left)(z2_right)] (zoneL2) {};
\end{scope}
\node[zonetitle=blkZoneL2, anchor=west] at ([xshift=0.32cm]zoneL2.north west)
    {Actuaci\'on y Motor (Cascada)};

% ==============================================================================
% ZONA 3: SENSADO BASE Y LPF IIR 1 (DERECHA -> IZQUIERDA)
% ==============================================================================
\coordinate (z3_y) at ([yshift=-1.65cm]zoneL2.south);
\node[sensTblk, anchor=center] (as5600_base) at (puente.center |- z3_y) {AS5600\\Base};
\node[sensTblk, left=\scSensGapX of as5600_base] (desenv) {Desenvoltura};
\node[sensTblk, left=\scSensGapX of desenv] (deriv_t) {Derivador $\theta$};
\node[sensTblk, left=\scSensGapX of deriv_t] (lpf1_t) {LPF IIR 1};

\coordinate (z3_top)   at ([yshift=0.42cm]lpf1_t.north);
\coordinate (z3_bot)   at ([yshift=-0.32cm]lpf1_t.south);
\coordinate (z3_left)  at ([xshift=-0.78cm]lpf1_t.west);
\coordinate (z3_right) at ([xshift=0.30cm]as5600_base.east);

\begin{scope}[on background layer]
    \node[zonebox=blkZoneL3, inner sep=0pt,
          fit=(z3_top)(z3_bot)(z3_left)(z3_right)] (zoneL3) {};
\end{scope}
\node[zonetitle=blkZoneL3, anchor=west] at ([xshift=0.32cm]zoneL3.south west)
    {Sensado Base y LPF IIR 1};

% ==============================================================================
% CONEXIONES INTERNAS ZONA 1
% ==============================================================================
\coordinate (trk_dt1_x)   at ([xshift=-0.52cm]zone1.west);
\coordinate (trk_theta_x) at ([xshift=-\scTrkStepX]trk_dt1_x);
\coordinate (ref_in)      at ([xshift=-1.35cm]trk_theta_x |- sig_p.west);

% Referencia theta_ref
\draw[sig] (ref_in) -- (sig_p.west);
\node[wirelbl, anchor=south] at ($(ref_in)!0.50!(trk_theta_x |- sig_p.west) + (0, 2pt)$) {$\theta_{\mathrm{ref}}$};
\node[signlbl, anchor=south] at ($(sig_p.west)+(-0.15, 0.08)$) {$+$};

% Error e_theta y bifurcación hacia P, I y Reinicio Clegg
\draw[wire] (sig_p.east) -- (split_e)
    node[wirelbl, midway, above=2pt] {$e_\theta$};
\node[dot=blkInk] at (split_e) {};

\draw[sigshort] (split_e) -- (win_box.west);
\draw[sigshort] (win_box.east) -- (int_box.west);
\draw[sigshort] (int_box.east) -- (clamp_box.west);
\draw[sigshort] (clamp_box.east) -- (sig_c.west)
    node[wirelbl, pos=0.38, above=2pt] {$u_i$};
\node[signlbl, anchor=south] at ($(sig_c.west)+(-0.15, 0.08)$) {$+$};

\coordinate (clegg_tap) at (split_e |- clegg_box.west);
\draw[wire] (split_e) -- (clegg_tap);
\node[dot=blkInk] at (clegg_tap) {};
\draw[sigshort] (clegg_tap) -- (clegg_box.west);
\draw[sig, rounded corners=4pt] (clegg_tap) |- (prop_box.west);

\draw[sigClegg] (clegg_box.south) -- (int_box.north)
    node[wirelbl, text=blkDistEdge, midway, right=3pt] {$\mathrm{reset}$};

\draw[sig, rounded corners=4pt] (prop_box.east) -| (sig_c.north)
    node[wirelbl, pos=0.75, right=3pt] {$u_p$};
\node[signlbl, anchor=south west] at ($(sig_c.north)+(0.08, 0.08)$) {$+$};

\draw[sig, rounded corners=4pt] (deriv_box.east) -| (sig_c.south)
    node[wirelbl, pos=0.25, below=2pt] {$u_d$};
\node[signlbl, anchor=north west] at ($(sig_c.south)+(0.08, -0.08)$) {$+$};

\draw[sigshort] (sig_c.east) -- (sat_box.west);
\draw[sigshort] (sat_box.east) -- (slew_box.west);
\draw[sigshort] (slew_box.east) -- (sel_box.west);

% Bus interno de velocidad dtheta_f1 -> Derivativo y Histéresis (con saltos limpios)
\coordinate (vel_cross_e) at (split_e |- deriv_box.west);
\coordinate (vel_tap)     at (win_box.center |- deriv_box.west);
\coordinate (vel_cross_r) at (int_box.center |- deriv_box.west);

\coordinate (ehop_top) at ($(vel_cross_e)+(0, \scHopRad)$);
\coordinate (ehop_bot) at ($(vel_cross_e)+(0,-\scHopRad)$);
\coordinate (rhop_top) at ($(vel_cross_r)+(0, \scHopRad)$);
\coordinate (rhop_bot) at ($(vel_cross_r)+(0,-\scHopRad)$);

\coordinate (hyst_in_e) at ([yshift=-\scHystInOff]hyst_box.west);
\coordinate (hyst_in_v) at ([yshift=+\scHystInOff]hyst_box.west);

% e_theta -> Histéresis (salto sobre dtheta_f1)
\draw[wire] (split_e) -- (ehop_top);
\draw[wire] (ehop_top) arc[start angle=90, end angle=-90, radius=\scHopRad];
\draw[sig, rounded corners=4pt] (ehop_bot) |- (hyst_in_e)
    node[wirelbl, pos=0.78, below=2pt] {$e_\theta$};

% reset desde Histéresis -> Integrador (salto sobre dtheta_f1)
\draw[wireBrakeDash] (hyst_box.north) -- (rhop_bot)
    node[wirelbl, text=blkActEdge, midway, right=3pt] {$\mathrm{reset}$};
\draw[wire, draw=blkActEdge] (rhop_bot) arc[start angle=-90, end angle=90, radius=\scHopRad];
\draw[sigBrakeDash] (rhop_top) -- (int_box.south);

% Histéresis -> Freno Activo y lock -> Selector Freno
\coordinate (brake_en_tap) at (brake_box.south |- hyst_box.east);
\coordinate (sel_cmd_in)   at ([xshift=+\scSelOffX]sel_box.south);
\coordinate (sel_brk_in)   at ([xshift=-\scSelOffX]sel_box.south);

\draw[wireBrakeDash] (hyst_box.east) -- (brake_en_tap);
\node[dot=blkActEdge] at (brake_en_tap) {};
\draw[sigBrakeDash] (brake_en_tap) -- (brake_box.south);
\draw[sigBrakeDash, rounded corners=4pt] (brake_en_tap) -| (sel_cmd_in)
    node[wirelbl, text=blkActEdge, pos=0.75, right=3pt] {$\mathrm{lock}$};
\draw[sigBrakeSolid, rounded corners=4pt] (brake_box.east) -| (sel_brk_in);

% ==============================================================================
% CONEXIONES ZONA 1 -> ZONA 2 (TENSIÓN u_servo Y DESVÍO FRENO ACTIVO)
% ==============================================================================
\coordinate (u_right_x)  at ([xshift=0.68cm]zone1.east);
\coordinate (lane_u_y)   at ([yshift=0.68cm]zoneL2.north);
\coordinate (u_in_x)     at ([xshift=-0.38cm]zoneL2.west);
\coordinate (sel_out_u)  at ([yshift= 0.14cm]sel_box.east);
\coordinate (sel_out_br) at ([yshift=-0.14cm]sel_box.east);

% Canal tensión u_servo -> Banda Silenciosa
\draw[sig, rounded corners=4pt]
    (sel_out_u) -- (u_right_x |- sel_out_u)
    -- (u_right_x |- lane_u_y)
    -- (u_in_x |- lane_u_y) |- (banda.west);
\node[wirelbl, anchor=south]
    at ($(rampa.center |- lane_u_y)!0.50!(pwm.center |- lane_u_y) + (0, 1.5pt)$) {$u_{\mathrm{servo}}$};

% Canal Freno directo -> Puente H L298N (en corredor limpio entre lane_u_y y zoneL2.north)
\coordinate (br_right_x) at ([xshift=0.32cm]zone1.east);
\coordinate (lane_br_y)  at ([yshift=0.30cm]zoneL2.north);
\draw[sigBrakeDash, rounded corners=4pt]
    (sel_out_br) -- (br_right_x |- sel_out_br)
    -- (br_right_x |- lane_br_y) -| (puente.north);
\node[wirelbl, text=blkActEdge, anchor=south]
    at ($(puente.east |- lane_br_y)!0.52!(motor.east |- lane_br_y) + (0, 1pt)$) {Freno};

% Cascada interna Zona 2
\draw[sigshort] (banda.east) -- (rampa.west);
\draw[sigshort] (rampa.east) -- (pwm.west);
\draw[sigshort] (pwm.east) -- (puente.west);
\draw[sigshort] (puente.east) -- node[wirelbl, midway, above=2pt]{$V_m$} (motor.west);

% ==============================================================================
% CONEXIONES ZONA 3 -> 2 PISTAS SEPARADAS (POSICIÓN theta Y VELOCIDAD dtheta_f1)
% ==============================================================================
\draw[fbT, rounded corners=4pt] (motor.south) |- (as5600_base.east);
\node[wirelbl, text=blkArmEdge, anchor=west]
    at ([xshift=0.08cm]$(motor.south |- zoneL2.south)!0.48!(motor.south |- zoneL3.north)$) {$\theta(t)$};

\draw[fbT] (as5600_base.west) -- (desenv.east);
\draw[fbT] (desenv.west) -- (deriv_t.east);
\draw[fbT] (deriv_t.west) -- (lpf1_t.east);

% 1) PISTA INTERIOR (VELOCIDAD dtheta_f1) -> Derivativo (-dtheta_f1) y Histéresis (dtheta_f1)
\draw[fbT, rounded corners=4pt]
    (lpf1_t.west) -- (trk_dt1_x |- lpf1_t.west)
    -- (trk_dt1_x |- deriv_box.west) -- (deriv_box.west);
\node[wirelbl, text=blkArmEdge, anchor=south]
    at ($(lpf1_t.west)!0.50!(zoneL3.west |- lpf1_t.west) + (0, 2pt)$) {$\dot{\theta}_{f1}$};
\node[wirelbl, text=blkArmEdge, anchor=south]
    at ($(sig_p.west |- deriv_box.west)!0.42!(vel_cross_e) + (0, 2pt)$) {$\dot{\theta}_{f1}$};
\node[wirelbl, text=blkArmEdge, anchor=south]
    at ($(vel_cross_r)!0.54!(deriv_box.west) + (0, 2pt)$) {$-\dot{\theta}_{f1}$};

\node[dot=blkArmEdge] at (vel_tap) {};
\draw[fbT, rounded corners=4pt] (vel_tap) |- (hyst_in_v)
    node[wirelbl, text=blkArmEdge, pos=0.78, above=2pt] {$\dot{\theta}_{f1}$};

% 2) PISTA EXTERIOR (POSICIÓN theta) -> Sigma_p (-) [salta sobre trk_dt1_x limpiamente]
\coordinate (tap_theta)    at ($(desenv.west)!0.56!(deriv_t.east)$);
\coordinate (rail_theta_y) at ($(zoneL2.south)!0.52!(zoneL3.north)$);
\coordinate (y_sigp_fb)    at ([yshift=-0.46cm]sig_p.south);

\node[dot=blkArmEdge] at (tap_theta) {};
\draw[lineT, rounded corners=4pt]
    (tap_theta) -- (tap_theta |- rail_theta_y)
    -- ($(trk_dt1_x |- rail_theta_y)+(\scHopRad, 0)$)
    arc[start angle=0, end angle=180, radius=\scHopRad];
\draw[fbT, rounded corners=4pt]
    ($(trk_dt1_x |- rail_theta_y)+(-\scHopRad, 0)$) -- (trk_theta_x |- rail_theta_y)
    -- (trk_theta_x |- y_sigp_fb) -| (sig_p.south);
\node[wirelbl, text=blkArmEdge, anchor=north]
    at ($(deriv_t.north |- rail_theta_y) + (-0.15cm, -1pt)$) {$\theta$};
\node[wirelbl, text=blkArmEdge, anchor=south]
    at ($(trk_dt1_x |- y_sigp_fb)!0.52!(sig_p.south |- y_sigp_fb) + (0, 1.5pt)$) {$\theta$};
\node[signlbl, text=blkArmEdge, anchor=north west] at ($(sig_p.south)+(0.08, -0.06)$) {$-$};

\end{tikzpicture}%
}
\caption{Diagrama de bloques at\'omico exclusivo del control servo a $\qty{0}{\degree}$ activo en \texttt{FURUTA\_STATE\_MOTOR\_PID\_POS}, \texttt{STOP\_HOME} y la subfase \texttt{F1} de \texttt{RECOVER\_HOME} (\texttt{Furuta\_DriveArmToZero\_1kHz} en \texttt{furuta\_control.c}, $T_s=\qty{1.0}{\milli\second}$). Combina tres ramas paralelas ---proporcional, integral no lineal con ventana de $\pm\qty{5}{\degree}$, sujeci\'on en $\pm\qty{6.5}{\volt}$ y reinicio instant\'aneo de Clegg~\eqref{eq:clegg_pos}, y derivativa directa sobre $-\dot{\theta}_{f1}$--- seguidas por saturaci\'on ($\pm\qty{5.8}{\volt}$), limitador de pendiente ($\qty{0.055}{\volt\per\milli\second}$), selector con freno activo \texttt{L298N\_Brake}, la cascada de potencia L298N hacia el motor $80:1$ y la realimentaci\'on base AS5600 con se\~nales separadas de posici\'on ($\theta$) y velocidad filtrada ($\dot{\theta}_{f1}$).}
\label{fig:detalle_servo_clegg}
\end{figure*}
